% 2011 年广东卷理科第 7 题：某几何体的三视图（正视图为平行四边形，侧视图、俯视图为矩形）。
% 来源 PDF：普通高考/2011/2011广东理.pdf，SHA-256 676f6eb7aaa62d199da64d67e0951301d980eb1e4d82193579afc4a4637ac874。
% 原页：第 1 页（A3 三栏）第 1 栏下部，书页码 663。
% 重绘依据（源为矢量，坐标取自 pdftocairo SVG，真页 pt、原点左上、y 向下）：
%   正视图平行四边形：下边 (124.68,637.96)-(201.18,637.96)，上边 (150.18,593.79)-(226.68,593.79)，
%     即水平边 3 单位、斜边 2 单位、竖向高 √3 单位（76.5pt = 3 单位，51.0pt = 2 单位，44.17pt = √3 单位）；
%   侧视图矩形：(252.18,593.79)-(328.69,637.96)，宽 3、高 √3；
%   俯视图矩形：(124.68,663.46)-(226.68,739.97)，宽 4、高 3，内部 x=150.18 处为实线、x=201.18 处为虚线，
%     下方与上方均标 1、2、1（即 30.18、81.18 两处分界）。
%   本图取本地坐标 (x,y)=(x_page−120, 760−y_page)，1 单位 = 1 TeX pt，y 向上；尺寸标注照原页补画，
%   拓扑与源一致，不按原页绝对比例。
\begin{tikzpicture}[x=1pt, y=1pt, >=stealth, semithick, font=\small]
  % ---- 正视图（平行四边形：水平边 3，斜边 2，高 √3）----
  \draw (4.68,122.04) -- (81.18,122.04) -- (106.68,166.21) -- (30.18,166.21) -- cycle;
  % ---- 侧视图（矩形 3×√3）----
  \draw (132.18,122.04) rectangle (208.69,166.21);
  % ---- 俯视图（矩形 4×3，内部实线与虚线各一）----
  \draw (4.68,12.03) rectangle (106.68,88.54);
  \draw (30.18,12.03) -- (30.18,88.54);
  \draw[dashed] (81.18,12.03) -- (81.18,88.54);
  % ---- 正视图尺寸：上边 3 ----
  \draw (30.18,168.6) -- (30.18,171.3);
  \draw (106.68,168.6) -- (106.68,171.3);
  \draw[<->] (30.18,171.3) -- (106.68,171.3) node[midway,above] {$3$};
  % ---- 正视图尺寸：斜边 2 ----
  \draw[<->] (0.68,124.34) -- (26.18,168.51) node[midway,above left,sloped,inner sep=1.5pt] {$2$};
  % ---- 侧视图尺寸：宽 3 ----
  \draw (132.18,168.6) -- (132.18,171.3);
  \draw (208.69,168.6) -- (208.69,171.3);
  \draw[<->] (132.18,171.3) -- (208.69,171.3) node[midway,above] {$3$};
  % ---- 俯视图尺寸：上沿 1、2、1 ----
  \draw[<->] (4.68,93.64) -- (30.18,93.64) node[midway,above] {$1$};
  \draw[<->] (30.18,93.64) -- (81.18,93.64) node[midway,above] {$2$};
  \draw[<->] (81.18,93.64) -- (106.68,93.64) node[midway,above] {$1$};
  % ---- 俯视图尺寸：下沿 1、2、1 ----
  \draw[<->] (4.68,3.0) -- (30.18,3.0) node[midway,below] {$1$};
  \draw[<->] (30.18,3.0) -- (81.18,3.0) node[midway,below] {$2$};
  \draw[<->] (81.18,3.0) -- (106.68,3.0) node[midway,below] {$1$};
  % ---- 视图名称 ----
  \node at (55.75,114.0) {正视图};
  \node at (170.55,114.0) {侧视图};
  \node at (55.75,-13.5) {俯视图};
\end{tikzpicture}
